% !TEX program = tectonic
% =============================================================================
%  AULA 11 — IA e Machine Learning nas Redes sem Fio
%  Curso: Redes sem Fio  |  Profa. Anelise Munaretto
%  ML para otimização, predição, alocação de recursos.
%  Compilar: tectonic aula11.tex
% =============================================================================
\documentclass[aspectratio=169,11pt]{beamer}
% =============================================================================
%  PREÁMBULO COMÚN para as aulas de Redes sem Fio (Beamer + TikZ)
%  Incluido com % =============================================================================
%  PREÁMBULO COMÚN para as aulas de Redes sem Fio (Beamer + TikZ)
%  Incluido com % =============================================================================
%  PREÁMBULO COMÚN para as aulas de Redes sem Fio (Beamer + TikZ)
%  Incluido com \input{preamble.tex} en cada aula.
%  Paleta de cores e estilos são definidos aqui uma única vez.
% =============================================================================

\usepackage[T1]{fontenc}
\usepackage{amsmath,amssymb,mathtools}
\usepackage{graphicx}
\usepackage{tikz}
\usetikzlibrary{positioning,arrows.meta,calc,backgrounds,decorations.pathmorphing,fit,shadows,shapes.symbols,shapes.geometric}
\usepackage{pgfplots}
\pgfplotsset{compat=1.16}
\usepackage{tabularx,booktabs,multirow,array}
\usepackage[normalem]{ulem}
\usepackage{hyperref}

% ---------- Paleta de cores própria ----------
\definecolor{aneliseAzul}{RGB}{20,60,110}
\definecolor{aneliseVerde}{RGB}{30,140,70}
\definecolor{aneliseLaranja}{RGB}{215,120,20}
\definecolor{aneliseGris}{RGB}{90,90,90}

% ---------- Tema ----------
\usetheme{Madrid}
\usecolortheme{whale}
\setbeamercolor{structure}{fg=aneliseAzul}
\setbeamercolor{title}{fg=aneliseAzul}
\setbeamercolor{frametitle}{fg=aneliseAzul}
\setbeamercolor{block title}{fg=aneliseAzul,bg=aneliseGris!12!white}
\setbeamercolor{block body}{bg=aneliseGris!7!white}
\hypersetup{colorlinks=true,linkcolor=aneliseAzul,urlcolor=aneliseVerde}

% ---------- Comandos auxiliares ----------
\newcommand{\kurose}{Kurose \& Ross, \emph{Computer Networking: a Top-Down Approach}}
\newcommand{\forouzan}{Forouzan, \emph{Data Communications and Networking}}
\newcommand{\stallings}{Stallings, \emph{Wireless Communications \& Networks}}
\newcommand{\tanenbaum}{Tanenbaum, \emph{Computer Networks}}
\newcommand{\rfc}[1]{\href{https://www.rfc-editor.org/rfc/rfc#1.txt}{RFC~#1}}
% -----------------------------------------------------------------------------
%  Estilos TikZ comuns para desenhar redes (posicionamento relativo)
% -----------------------------------------------------------------------------
\tikzset{
  host/.style={draw,rounded corners=2mm,fill=aneliseAzul!15,inner sep=1.5mm,font=\scriptsize},
  rt/.style={draw,rounded corners=1mm,fill=aneliseLaranja!25,inner sep=1mm,font=\scriptsize},
  nube/.style={draw,cloud,aspect=2.2,fill=aneliseGris!15,inner sep=1.5mm,font=\scriptsize},
  el/.style={draw,rounded corners=1.2mm,fill=aneliseGris!18,inner sep=0.8mm,font=\scriptsize},
  enl/.style={thick},
  enlA/.style={thick,-Stealth},
} en cada aula.
%  Paleta de cores e estilos são definidos aqui uma única vez.
% =============================================================================

\usepackage[T1]{fontenc}
\usepackage{amsmath,amssymb,mathtools}
\usepackage{graphicx}
\usepackage{tikz}
\usetikzlibrary{positioning,arrows.meta,calc,backgrounds,decorations.pathmorphing,fit,shadows,shapes.symbols,shapes.geometric}
\usepackage{pgfplots}
\pgfplotsset{compat=1.16}
\usepackage{tabularx,booktabs,multirow,array}
\usepackage[normalem]{ulem}
\usepackage{hyperref}

% ---------- Paleta de cores própria ----------
\definecolor{aneliseAzul}{RGB}{20,60,110}
\definecolor{aneliseVerde}{RGB}{30,140,70}
\definecolor{aneliseLaranja}{RGB}{215,120,20}
\definecolor{aneliseGris}{RGB}{90,90,90}

% ---------- Tema ----------
\usetheme{Madrid}
\usecolortheme{whale}
\setbeamercolor{structure}{fg=aneliseAzul}
\setbeamercolor{title}{fg=aneliseAzul}
\setbeamercolor{frametitle}{fg=aneliseAzul}
\setbeamercolor{block title}{fg=aneliseAzul,bg=aneliseGris!12!white}
\setbeamercolor{block body}{bg=aneliseGris!7!white}
\hypersetup{colorlinks=true,linkcolor=aneliseAzul,urlcolor=aneliseVerde}

% ---------- Comandos auxiliares ----------
\newcommand{\kurose}{Kurose \& Ross, \emph{Computer Networking: a Top-Down Approach}}
\newcommand{\forouzan}{Forouzan, \emph{Data Communications and Networking}}
\newcommand{\stallings}{Stallings, \emph{Wireless Communications \& Networks}}
\newcommand{\tanenbaum}{Tanenbaum, \emph{Computer Networks}}
\newcommand{\rfc}[1]{\href{https://www.rfc-editor.org/rfc/rfc#1.txt}{RFC~#1}}
% -----------------------------------------------------------------------------
%  Estilos TikZ comuns para desenhar redes (posicionamento relativo)
% -----------------------------------------------------------------------------
\tikzset{
  host/.style={draw,rounded corners=2mm,fill=aneliseAzul!15,inner sep=1.5mm,font=\scriptsize},
  rt/.style={draw,rounded corners=1mm,fill=aneliseLaranja!25,inner sep=1mm,font=\scriptsize},
  nube/.style={draw,cloud,aspect=2.2,fill=aneliseGris!15,inner sep=1.5mm,font=\scriptsize},
  el/.style={draw,rounded corners=1.2mm,fill=aneliseGris!18,inner sep=0.8mm,font=\scriptsize},
  enl/.style={thick},
  enlA/.style={thick,-Stealth},
} en cada aula.
%  Paleta de cores e estilos são definidos aqui uma única vez.
% =============================================================================

\usepackage[T1]{fontenc}
\usepackage{amsmath,amssymb,mathtools}
\usepackage{graphicx}
\usepackage{tikz}
\usetikzlibrary{positioning,arrows.meta,calc,backgrounds,decorations.pathmorphing,fit,shadows,shapes.symbols,shapes.geometric}
\usepackage{pgfplots}
\pgfplotsset{compat=1.16}
\usepackage{tabularx,booktabs,multirow,array}
\usepackage[normalem]{ulem}
\usepackage{hyperref}

% ---------- Paleta de cores própria ----------
\definecolor{aneliseAzul}{RGB}{20,60,110}
\definecolor{aneliseVerde}{RGB}{30,140,70}
\definecolor{aneliseLaranja}{RGB}{215,120,20}
\definecolor{aneliseGris}{RGB}{90,90,90}

% ---------- Tema ----------
\usetheme{Madrid}
\usecolortheme{whale}
\setbeamercolor{structure}{fg=aneliseAzul}
\setbeamercolor{title}{fg=aneliseAzul}
\setbeamercolor{frametitle}{fg=aneliseAzul}
\setbeamercolor{block title}{fg=aneliseAzul,bg=aneliseGris!12!white}
\setbeamercolor{block body}{bg=aneliseGris!7!white}
\hypersetup{colorlinks=true,linkcolor=aneliseAzul,urlcolor=aneliseVerde}

% ---------- Comandos auxiliares ----------
\newcommand{\kurose}{Kurose \& Ross, \emph{Computer Networking: a Top-Down Approach}}
\newcommand{\forouzan}{Forouzan, \emph{Data Communications and Networking}}
\newcommand{\stallings}{Stallings, \emph{Wireless Communications \& Networks}}
\newcommand{\tanenbaum}{Tanenbaum, \emph{Computer Networks}}
\newcommand{\rfc}[1]{\href{https://www.rfc-editor.org/rfc/rfc#1.txt}{RFC~#1}}
% -----------------------------------------------------------------------------
%  Estilos TikZ comuns para desenhar redes (posicionamento relativo)
% -----------------------------------------------------------------------------
\tikzset{
  host/.style={draw,rounded corners=2mm,fill=aneliseAzul!15,inner sep=1.5mm,font=\scriptsize},
  rt/.style={draw,rounded corners=1mm,fill=aneliseLaranja!25,inner sep=1mm,font=\scriptsize},
  nube/.style={draw,cloud,aspect=2.2,fill=aneliseGris!15,inner sep=1.5mm,font=\scriptsize},
  el/.style={draw,rounded corners=1.2mm,fill=aneliseGris!18,inner sep=0.8mm,font=\scriptsize},
  enl/.style={thick},
  enlA/.style={thick,-Stealth},
}

\title[Redes sem Fio]{Redes sem Fio\\ IA e Machine Learning\\ nas Redes sem Fio}
\subtitle{Aula 11}
\author[Anelise Munaretto]{Profa. Anelise Munaretto}
\institute[Curso]{Redes sem Fio --- Ciência da Computação}
\date{2026}

\begin{document}

\begin{frame}[plain]
  \titlepage
\end{frame}

\begin{frame}{Objetivos de aprendizagem}
  Ao final desta aula, você será capaz de:
  \begin{itemize}
    \item \textbf{descrever} o ciclo agente--ambiente do RL e \textbf{diferenciar} aprendizado centralizado de \emph{federated learning};
    \item \textbf{distinguir} aprendizagem supervisionada, não-supervisionada, por reforço e profunda;
    \item \textbf{aplicar} a equação de Bellman e a atualização Q-learning em alocação de recursos;
    \item \textbf{explicar} o uso de LSTM na predição de mobilidade e \emph{handover};
    \item \textbf{avaliar} desafios de dados, privacidade, confiabilidade e padronização de ML em redes (3GPP RM-NET, O-RAN).
  \end{itemize}
  \begin{alertblock}{Pré-requisitos}
    Aulas 7 (5G e redes móveis) e 10 (drones, satélites e NTN).
  \end{alertblock}
\end{frame}

\section{Motivação}

\begin{frame}{Por que IA/ML em redes sem fio?}
  \begin{block}{Crise de complexidade}
    Redes 5G/6G têm centenas de parâmetros ajustáveis; otimização manual é inviável. IA/ML oferece soluções para:
  \end{block}
  \begin{itemize}
    \item predição de tráfego e mobilidade
    \item alocação dinâmica de recursos (espectro, potência, beam)
    \item detecção de anomalias e falhas
    \item gerenciamento autônomo (\emph{zero-touch})
    \item otimização cross-layer
  \end{itemize}
\end{frame}

\section{Técnicas}

\begin{frame}{Técnicas de ML aplicadas}
  \begin{columns}
    \begin{column}{0.48\textwidth}
      \textbf{Aprendizagem Supervisionada}
      \begin{itemize}
        \item classificação de interferência
        \item predição de SNR/CQI
        \item detecção de tipo de tráfego
      \end{itemize}
      \textbf{Aprendizagem Não-Supervisionada}
      \begin{itemize}
        \item agrupamento de padrões de mobilidade
        \item compressão de CSI
      \end{itemize}
    \end{column}
    \begin{column}{0.48\textwidth}
      \textbf{Aprendizagem por Reforço (RL)}
      \begin{itemize}
        \item \emph{resource block allocation}
        \item escalonamento de pacotes
        \item controle de admissão
      \end{itemize}
      \textbf{Aprendizagem Profunda (DL)}
      \begin{itemize}
        \item \emph{channel estimation} com autoencoders
        \item beamforming com CNN
        \item predição de séries temporais (LSTM)
      \end{itemize}
    \end{column}
  \end{columns}
\end{frame}

\begin{frame}{RL para alocação de recursos}
  \begin{center}
    \begin{tikzpicture}[font=\scriptsize,
        c/.style={draw,rounded corners=2mm,fill=aneliseAzul!13,inner sep=1mm,text width=2.8cm,align=center},
      ]
      \node[c](agent){Agente\\ (alocador)};
      \node[c,right=2.4cm of agent](env){Ambiente\\ (rede sem fio)};
      \draw[->,thick,aneliseLaranja] (agent) -- node[midway,above=1mm,font=\tiny]{ação $A_t$ (alocação de RB)} (env);
      \draw[->,thick,aneliseVerde] (env.south) to[bend right=20] node[midway,below=1mm,font=\tiny]{$R_{t+1}$ (recompensa, ex.: throughput)} (agent.south);
      \draw[->,thick,aneliseGris] (env.north) to[bend left=20] node[midway,above=1mm,font=\tiny]{$S_{t+1}$ (novo estado)} (agent.north);
    \end{tikzpicture}
    \\[1mm] {\scriptsize \textbf{Figura 1:} ciclo agente--ambiente do aprendizado por reforço aplicado à alocação de recursos.}
  \end{center}
  \begin{itemize}
    \item Exemplo: alocação de RBs em 5G usando Deep Q-Learning (DQN)
    \item Estado: filas CQI, buffer; Ação: RBs por usuário; Recompensa: throughput $\times$ justiça (Jain)
  \end{itemize}
  \note{Enfatizar que o agente aprende por tentativa e erro maximizando a recompensa acumulada descontada ($G_t = \sum \gamma^k R_{t+k}$); a rede sem fio é o ambiente que transita de estado ao aplicar a ação de alocação.}
\end{frame}

\begin{frame}{Exemplo prático: predição de mobilidade}
  \begin{itemize}
    \item \textbf{Problema}: antecipar handover para evitar perda de pacotes
    \item \textbf{Método}: LSTM treinado com sequências de RSSI/RSRP
    \item \textbf{Entrada}: série temporal de 10 medições de potência (5 vizinhos)
    \item \textbf{Saída}: célula-alvo e tempo estimado de handover
    \item \textbf{Resultado}: redução de 40\% em falhas de handover (trabalhos reais)
  \end{itemize}
  \begin{block}{Implementação}
    Python + TensorFlow/Keras, treino offline com dados de drive-test, inferência no MEC ou no próprio UE.
  \end{block}
  \note{A LSTM processa janelas de 10 medições de RSSI/RSRP e prevê a célula-alvo do handover; destacar que o ganho real relatado na literatura é da ordem de 30--40\% de redução de falhas.}
\end{frame}

\section{Desafios}

\begin{frame}{Desafios de ML em redes sem fio}
  \begin{columns}
    \begin{column}{0.5\textwidth}
      \textbf{Dados}
      \begin{itemize}
        \item coleta de dados rotulados (CSI, tráfego)
        \item privacidade (GDPR, nuvem vs.\ borda)
        \item distribuição não-estacionária (mobilidade)
      \end{itemize}
      \textbf{Complexidade}
      \begin{itemize}
        \item inferência em tempo real (latência)
        \item tamanho do modelo (memória do dispositivo)
        \item \emph{federated learning} para escalar
      \end{itemize}
    \end{column}
    \begin{column}{0.5\textwidth}
      \textbf{Confiabilidade}
      \begin{itemize}
        \item modelos precisam generalizar
        \item interpretabilidade (explainable AI)
        \item robustez a ataques adversários
      \end{itemize}
      \textbf{Padronização}
      \begin{itemize}
        \item 3GPP Rel.\ 18: RM-NET (ML para RAN)
        \item O-RAN Alliance: rApp e xApp com ML
        \item ITU-T ML5G
      \end{itemize}
    \end{column}
  \end{columns}
  \note{Federated learning resolve privacidade e escala: o modelo é treinado localmente em cada gNB/UE e apenas os gradientes são agregados; correlacionar com a padronização O-RAN (rApp/xApp com ML).}
\end{frame}

\begin{frame}{Exercício resolvido --- Bellman e atualização Q-learning}
  \textbf{Problema:} um agente RL decide a alocação de um RB. A ação atual não gera recompensa imediata ($r = 0$), mas leva a um estado em que a melhor ação futura renderá $r = 10$ (estado terminal). Considere $\gamma = 0{,}9$ e partida com $Q(s,a) = 0$.
  \begin{enumerate}
    \item \textbf{Valor-alvo (Bellman):}
      $$Q(s,a) = r + \gamma\,\max_{a'} Q(s',a') = 0 + 0{,}9 \times 10 = 9.$$
    \item \textbf{Atualização Q-learning} com $\alpha = 0{,}1$:
      $$Q(s,a) \leftarrow Q(s,a) + \alpha\left[ r + \gamma\,\max_{a'} Q(s',a') - Q(s,a) \right] = 0 + 0{,}1\,(9 - 0) = 0{,}9.$$
    \item \textbf{Convergência:} repetindo a mesma experiência, $Q(s,a)$ se aproxima assintoticamente de $9$.
  \end{enumerate}
  \begin{alertblock}{Interpretação}
    O fator $\gamma = 0{,}9$ desconta o valor futuro: uma recompensa de $10$ a um passo de distância vale $0{,}9 \times 10 = 9$ hoje.
  \end{alertblock}
\end{frame}

\begin{frame}{Resumo da aula}
  \begin{itemize}
    \item \textbf{ML em redes}: predição de tráfego/mobilidade, alocação dinâmica de espectro/potência/beam, detecção de anomalias, \emph{zero-touch}.
    \item \textbf{Paradigmas}: supervisionado (classificação), não-supervisionado (agrupamento), por reforço (alocação) e profundo (LSTM, CNN, autoencoders).
    \item \textbf{RL/DQN}: agente aloca RBs com base em estado (CQI, buffer) e recompensa (throughput $\times$ justiça de Jain); DQN usa redes neurais como aproximador de $Q$.
    \item \textbf{LSTM}: predição de \emph{handover} a partir de séries de RSSI/RSRP (redução de $\approx 40\%$ de falhas).
    \item \textbf{Desafios}: dados rotulados, privacidade (GDPR), não-estacionariedade, latência de inferência e robustez.
    \item \textbf{Padronização}: 3GPP Rel.~18 RM-NET, O-RAN (rApp/xApp) e ITU-T ML5G; \emph{federated learning} para escalar com privacidade.
  \end{itemize}
\end{frame}

\begin{frame}{Autoavaliação}
  \begin{enumerate}
    \item No aprendizado por reforço, o agente obtém do ambiente:
      \begin{itemize}
        \item (a) apenas o rótulo correto da ação
        \item (b) apenas a ação executada
        \item (c) o novo estado e a recompensa
      \end{itemize}
      \pause
      \textbf{Gabarito:} (c)

    \item Uma LSTM é mais adequada para:
      \begin{itemize}
        \item (a) predição de séries temporais, como RSSI/RSRP
        \item (b) ordenação de pacotes em filas
        \item (c) compressão de imagens de beamforming
      \end{itemize}
      \pause
      \textbf{Gabarito:} (a)

    \item A principal diferença do \emph{federated learning} para o treino centralizado é:
      \begin{itemize}
        \item (a) usar apenas redes neurais profundas
        \item (b) treinar com dados locais e agregar apenas os modelos/gradientes
        \item (c) dispensar a camada de rede física
      \end{itemize}
      \pause
      \textbf{Gabarito:} (b)

    \item Com $\gamma = 0{,}9$, $r = 0$ e $\max Q(s',a') = 10$, a equação de Bellman resulta em $Q(s,a) =$
      \begin{itemize}
        \item (a) $9$
        \item (b) $10$
        \item (c) $0{,}9$
      \end{itemize}
      \pause
      \textbf{Gabarito:} (a)
  \end{enumerate}
\end{frame}

\begin{frame}{Laboratório sugerido}
  \begin{block}{Q-learning em Python/gymnasium}
    \begin{itemize}
      \item Implementar Q-learning para alocação de canal/RB em um ambiente simplificado (ex.: \texttt{FrozenLake} ou um \texttt{env} customizado).
      \item Variar $\alpha$ e $\gamma$ e observar a curva de recompensa acumulada e a convergência da tabela $Q$.
      \item Extensão: substituir a tabela por uma rede neural (DQN) usando PyTorch/TensorFlow.
    \end{itemize}
  \end{block}
  \begin{block}{Demonstração de \emph{federated learning}}
    \begin{itemize}
      \item Usar \texttt{tensorflow-federated} (TFF) para treinar um classificador com clientes simulados.
      \item Comparar precisão e custo de comunicação entre treino centralizado e federado.
    \end{itemize}
  \end{block}
\end{frame}

\section{Bibliografia}

\begin{frame}{Bibliografia}
  \begin{itemize}
    \item C. Zhang et al., ``Deep learning in mobile and wireless networking: A survey'', \emph{IEEE Commun. Surv. Tut.}, 2019
    \item N. C. Luong et al., ``Applications of deep reinforcement learning in communications and networking'', \emph{IEEE Commun. Mag.}, 2019
    \item 3GPP TR 38.817, \emph{Study on RAN-centric data collection and utilization (RAN-CU)}
    \item O-RAN Alliance, \emph{O-RAN Working Group 2: AI/ML Workflow Description}
    \item ITU-T Y.3172, \emph{Architectural framework for ML in future networks}
  \end{itemize}
\end{frame}

\begin{frame}{Leitura recomendada}
  \begin{itemize}
    \item R. Sutton \& A. Barto, \emph{Reinforcement Learning: An Introduction}, MIT Press --- referência clássica para RL, Bellman e Q-learning.
    \item C. Zhang et al., ``Deep learning in mobile and wireless networking: A survey'', \emph{IEEE Commun. Surv. Tut.}, 2019 --- panorama geral de DL aplicado a redes.
    \item N. C. Luong et al., ``Applications of deep reinforcement learning in communications and networking'', \emph{IEEE Commun. Mag.}, 2019 --- DQN aplicado a alocação de recursos.
    \item O-RAN Alliance, \emph{WG2: AI/ML Workflow Description} --- como o ML entra na RAN (rApp/xApp).
    \item 3GPP TR 38.817 (RAN-CU) --- coleta e uso de dados centrados na RAN para ML.
  \end{itemize}
  \begin{alertblock}{Dica de estudo}
    Reimplemente o exercício de Bellman em Python e depois generalize para DQN; isso consolida o ciclo agente--ambiente.
  \end{alertblock}
\end{frame}

\end{document}